\documentclass{article}

% Recommended, but optional, packages for figures and better typesetting:
\usepackage{microtype}
\usepackage{graphicx}
\usepackage{subfigure}
\usepackage{booktabs} % for professional tables

% hyperref makes hyperlinks in the resulting PDF.
% If your build breaks (sometimes temporarily if a hyperlink spans a page)
% please comment out the following usepackage line and replace
% \usepackage{icml2021} with \usepackage[nohyperref]{icml2021} above.
\usepackage{hyperref}

% Attempt to make hyperref and algorithmic work together better:
\newcommand{\theHalgorithm}{\arabic{algorithm}}

% Use the following line for the initial blind version submitted for review:
%\usepackage{icml2021}

% If accepted, instead use the following line for the camera-ready submission:
\usepackage[accepted]{icml2021}

\begin{document}

\twocolumn[
\icmltitle{Your Full Project Title Here}

\begin{icmlauthorlist}
\icmlauthor{Author One}{}
\icmlauthor{Author Two}{}
\icmlauthor{Author Three}{}
\icmlauthor{Author Four}{}
\end{icmlauthorlist}

\vskip 0.3in
]

\section{Contributions}

Our project compares linear, regularized, and ensemble methods on
the Kaggle House Prices dataset. We started from a plain linear
regression and ended with a stacked ensemble of Lasso, Gradient
Boosting, and XGBoost. We did not chase the lowest RMSE only.
A lot of time went into figuring out why each model failed before
we moved to the next one. The OLS baseline, for example, gave
huge weights to sparse one-hot features. Once we understood that,
the move to Lasso and Ridge made sense, and the same kind of
reasoning led us to the tree-based models later on.

The full list of models we built and tuned is: OLS as a baseline;
Lasso and Ridge with cross-validated regularization; Gradient
Boosting with both grid search and randomized search; XGBoost; a
weighted hybrid of Lasso and Gradient Boosting; and three stacking
setups, with the final one being Lasso + Gradient Boosting +
XGBoost. We also compared raw features, StandardScaler, and the
Yeo--Johnson transform, and ran validation curves on six gradient
boosting hyperparameters to see how stable the tuned configuration
really was. A lot of the work in the Process section
(the Condition2\_PosN outlier diagnosis, the leaf/split stress
testing, etc.) does not show up in the final RMSE table but took
significant time, so we want to flag it here.

We used the Kaggle \emph{House Prices: Advanced Regression
Techniques} dataset and did not collect anything ourselves. The
link to the dataset and our code is in the
repository.\footnote{Dataset: \url{https://www.kaggle.com/competitions/house-prices-advanced-regression-techniques}.
Code: \url{https://github.com/jwguanusc/CSCI567-FINAL}.} The
repository has a single notebook plus a short list of commands in
the README to reproduce all the numbers in the report.

For libraries, scikit-learn did most of the heavy lifting
(cross-validation, hyperparameter search, the stacking and voting
regressors). XGBoost was used for the boosting model, and pandas,
NumPy, SciPy, Matplotlib, and Seaborn for the rest. Two public
Kaggle notebooks (jesucristo, 2017; mviola, 2018) helped us with
some of the feature-engineering choices.

\section{AI Usage}

We used AI tools (mostly Claude) during the project. How much we
relied on them depended on what we were doing.

We did not use AI much for ideation. The general direction --
comparing linear, regularized, and ensemble models on a tabular
regression task -- came from the course content and from looking
at top Kaggle solutions. The choice of which models to add at
which stage was ours. We sometimes asked AI for a quick second
opinion (for example, whether stacking on top of a strong gradient
boosting model was worth trying), but we did not let it pick the
direction.

Coding was where AI helped the most. A lot of the
\texttt{scikit-learn} setup is repetitive -- KFold splits,
\texttt{GridSearchCV} parameter grids, stacking pipelines, and
matplotlib plots like the validation curves and feature importance
charts. AI was a fast way to get the first version of those, which
we then read through, ran, and tweaked to fit our pipeline. We
also used it to figure out a few \texttt{scikit-learn} error
messages and parameter combinations that did not behave as we
expected.

For literature review we did not use AI at all. The references in
this paper were found through the course materials, the Kaggle
forum, and following citations from one paper to another, and we
read them ourselves. We were aware of the issue with hallucinated
citations, so we manually checked each reference against the
actual source.

Sometimes we used AI as a quick reference when a concept was not
clear, for example why \texttt{StandardScaler} alone is not enough
when the features are heavily skewed, or how out-of-fold
predictions feed into a stacking ensemble. These were starting
points -- the actual choices in our pipeline came from our own
experiments.

For writing, AI was used to polish the language of paragraphs we
had already drafted, and to fix awkward phrasing. The numbers, the
figures, and the analysis itself are ours, and we take
responsibility for the contents of this report.

\end{document}
